\documentclass[a4paper,12pt]{article}
\usepackage{fontspec}
\setmainfont{FreeSerif}
\setsansfont{FreeSans}
\setmonofont{FreeMono}
\usepackage{polyglossia}
\setmainlanguage{russian}
\usepackage[left=3cm,right=1.5cm,top=2cm,bottom=2cm]{geometry}
\usepackage{setspace}
\usepackage{enumitem}
\usepackage{longtable,array,booktabs,needspace}
\usepackage{tikz}
\usetikzlibrary{arrows.meta,positioning,calc,shapes.geometric}
\usepackage{graphicx,float}
\usepackage[unicode,hidelinks]{hyperref}

\setstretch{1.15}
\tolerance=3000
\emergencystretch=3em
\hbadness=10000
\setlength{\parindent}{1.25cm}
\setlength{\parskip}{0pt}
\setlist{itemsep=2pt,topsep=3pt}
\renewcommand{\arraystretch}{1.2}

% ---------- макросы ----------
\newenvironment{ucd}[3]{%
  \needspace{10\baselineskip}
  \begin{longtable}{|>{\bfseries\raggedright\arraybackslash}p{3.0cm}|p{11.8cm}|}
  \hline
  \multicolumn{2}{|l|}{\textbf{#1.\ #2}}\\ \hline
  Актор & #3\\ \hline
}{\end{longtable}}
\newcommand{\steps}[1]{\begin{enumerate}[nosep,leftmargin=*]#1\end{enumerate}}

\tikzset{
  uc/.style={ellipse,draw,minimum height=1.2cm,text width=4.9cm,align=center,font=\footnotesize\linespread{1}\selectfont,inner sep=0pt},
  box/.style={rectangle,draw,rounded corners=2pt,align=center,font=\small,minimum height=0.9cm},
  >=Stealth
}
\usetikzlibrary{positioning,calc}
\newcommand{\ent}[2]{\begin{tabular}{@{}p{3.4cm}@{}}\textbf{\hfil #1\hfil}\\\hline #2\end{tabular}}
\newcommand{\actor}[3]{% id, (x,y), подпись
  \begin{scope}[shift={#2}]
    \draw (0,0.55) circle (0.13);
    \draw (0,0.42)--(0,-0.05);
    \draw (-0.22,0.28)--(0.22,0.28);
    \draw (0,-0.05)--(-0.17,-0.4);
    \draw (0,-0.05)--(0.17,-0.4);
    \node[minimum size=0.9cm] (#1) at (0,0.1) {};
    \node[below,font=\small,align=center] at (0,-0.4) {#3};
  \end{scope}}

\begin{document}

% ================= ТИТУЛЬНЫЙ ЛИСТ =================
\begin{titlepage}
\centering
\singlespacing
Министерство науки и высшего образования Российской Федерации\\[2pt]
ФЕДЕРАЛЬНОЕ ГОСУДАРСТВЕННОЕ АВТОНОМНОЕ ОБРАЗОВАТЕЛЬНОЕ УЧРЕЖДЕНИЕ ВЫСШЕГО ОБРАЗОВАНИЯ\\[2pt]
«НАЦИОНАЛЬНЫЙ ИССЛЕДОВАТЕЛЬСКИЙ УНИВЕРСИТЕТ ИТМО»

\vspace{2.5cm}
Группа P3311

\vspace{1cm}
Дисциплина: Информационные системы

\vspace{3cm}
{\Large\textbf{Курсовая работа. Этап 1}}\\[8pt]
{\large Информационная система поиска научного руководителя\\
и подачи заявок на выпускную квалификационную работу}

\vspace{1cm}
Анализ предметной области, требования, прецеденты и архитектура

\vfill
\hfill\begin{minipage}{7.5cm}
\raggedleft
\textbf{Выполнили:}\\
Рязанов Григорий Алексеевич\\
Абрамова Анастасия Сергеевна\\[10pt]
\textbf{Преподаватель:}\\
Тюрин Иван Николаевич
\end{minipage}

\vspace{2cm}
Санкт-Петербург, 2026
\end{titlepage}

\tableofcontents
\newpage

% ================= 1 =================
\section{Предметная область}

\subsection{Выбор предметной области}
Предметная область — \textbf{выбор научного руководителя и закрепление за ним студента для выполнения выпускной квалификационной работы (ВКР, диплома)} на уровнях бакалавриата и магистратуры. Предметная область предлагается к согласованию с преподавателем (см.\ раздел~\ref{sec:agree}).

\subsection{Описание предметной области}

\paragraph{Участники.}
\begin{itemize}
  \item \textbf{Студент} (бакалавр или магистр) — выбирает руководителя и тему диплома. Одновременно у студента может быть только одна ВКР.
  \item \textbf{Преподаватель (научный руководитель)} — ведёт ограниченное число студентов, формулирует и публикует темы, отбирает студентов.
  \item \textbf{Администратор (деканат, учебный отдел)} — подтверждает подлинность преподавателей, рассматривает жалобы на отзывы и спорные статусы.
  \item \textbf{Внешняя система my.itmo} — источник данных об успеваемости студентов и статусе сдачи работы. Интеграция планируется после получения разрешения; до этого используется заглушка с тестовыми данными.
\end{itemize}

\paragraph{Основные понятия.}
\begin{description}[leftmargin=2cm,style=nextline]
  \item[Тема] — направление работы, которое преподаватель вёл или готов вести: название, описание, ключевые слова. В профиле преподавателя видны как текущие темы, так и темы, по которым один из его студентов уже сдал работу.
  \item[Заявка] — обращение студента к конкретному преподавателю. В заявке студент своими словами формулирует желаемую тему. Заявка имеет статус и положение в личном рейтинге студента.
  \item[Топ студента] — упорядоченный список из не более чем пяти заявок, который студент составляет для себя. Преподаватель положение заявки в топе не видит, и оно не влияет на шансы. Топ нужен, чтобы студент понимал, к кому он хочет больше, и какие там шансы.
  \item[Лимит мест] — заявленное преподавателем число студентов. Лимит «мягкий»: преподаватель может взять и больше (например, 22 при лимите 20), если не включил запрет приёма сверх нормы. В заполненность входят только принятые студенты.
  \item[Избранное] — отметка («звёздочка») преподавателя на заявке студента, которого он, скорее всего, возьмёт. Отметка не занимает место в лимите, но учитывается при оценке шансов других студентов, и все видят, сколько студентов «почти приняты».
  \item[Шанс] — оценка вероятности попасть к преподавателю в виде категории (низкий, средний, высокий).
  \item[Дата окончания отбора] — срок, который каждый преподаватель задаёт сам.
  \item[Отзыв] — оценки студента о руководителе после сдачи работы (по шкалам, анонимно) и необязательный текстовый комментарий (публичный, с именем автора).
\end{description}

\paragraph{Текущий процесс (как есть).}
Студенты ищут руководителя через неформальные каналы: списки кафедры, личные сообщения, рассказы однокурсников. У студента нет сведений о том, чем руководитель занимается на практике, насколько он строг и вовлечён, и каковы шансы попасть именно к нему. Преподаватель получает обращения вручную и до личного разговора не видит ни успеваемости, ни интересов студента. Деканат не имеет инструмента проверки достоверности данных. Итог: выбор делается наугад или по знакомству, часть студентов остаётся без руководителя к концу срока, а преподаватели получают нерелевантные обращения.

\paragraph{Целевой процесс (как будет).}
\begin{enumerate}
  \item Преподаватель регистрируется, указывает контакты, темы, лимит мест и дату окончания отбора. Деканат подтверждает его аккаунт.
  \item Студент ищет преподавателей по направлению и ключевым словам, изучает профили (темы, статистику отзывов, шанс) и подаёт до пяти заявок, формулируя тему сам.
  \item Студент может запросить контакты преподавателя, указанные при регистрации, и обсуждать работу вне системы.
  \item Преподаватель видит профиль и оценки студента и решает: отметить избранным, принять или отклонить (с причиной).
  \item Принятый студент подтверждает выбор, остальные его заявки отзываются. К дате окончания отбора система напоминает студенту и при отсутствии реакции принимает решение по правилам из п.~\ref{sec:rules}.
  \item После сдачи работы (статус из my.itmo) студенту предлагается оставить отзыв, статистика и сдача отражаются в профиле преподавателя.
\end{enumerate}
Сервис \textbf{не распределяет пары автоматически}: решение о сотрудничестве остаётся за студентом и преподавателем. Система снижает издержки поиска и даёт обеим сторонам данные для решения.

\subsection{Правила предметной области}\label{sec:rules}
\begin{enumerate}[label=R\arabic*.]
  \item У студента не более одной ВКР. Подать заявки можно только при отсутствии подтверждённой работы.
  \item Одновременно у студента не более 5 активных заявок. Положение заявки в топе видно только студенту.
  \item Преподаватель может отклонить заявку только с указанием причины.
  \item Заполненность лимита считается по принятым студентам (статусы «Принята» и «В работе»). Избранное в лимит не входит.
  \item Если лимит исчерпан или превышен, шанс для новых заявителей — «низкий», с пояснением, что при симпатии преподавателя можно попасть сверх нормы. Если преподаватель включил запрет приёма сверх нормы, приём заявок закрывается сразу по достижении лимита.
  \item Принятие считается состоявшимся после подтверждения студентом. При подтверждении все остальные заявки студента отзываются.
  \item За 24 часа до даты окончания отбора система уведомляет студентов, которых принял преподаватель, но которые ещё не подтвердили выбор. Если студент не отреагировал до даты окончания: заявка, стоящая в топе студента выше всех остальных, подтверждается автоматически; иначе отзывается.
  \item При смене руководителя студент указывает причину, заявка закрывается, место освобождается, все прошлые заявки студента сбрасываются.
  \item Оставить отзыв может только аккаунт, получивший статус «работа сдана». Оценки по шкалам анонимны (автора знает только система, это защита от накрутки). Текстовый отзыв необязателен и публикуется с именем.
  \item Шанс учитывает число заявок, число избранных, средний балл студента среди заявителей данного преподавателя и лимит с числом принятых. Преподаватель видит все оценки студента, но шанс считается только по среднему баллу.
  \item Если работа в my.itmo не сдана, итоговая оценка считается равной~2.
\end{enumerate}

% ================= 2 =================
\section{Назначение системы}

\subsection{Зачем нужна система}
Выбор руководителя ВКР — единственный этап обучения, где результат зависит от взаимного согласия двух сторон при ограниченных местах и отсутствии общих инструментов поиска. Информационная система вводит единое прозрачное пространство, в котором студент видит доступных преподавателей и свои шансы, а преподаватель получает структурированные заявки с данными об успеваемости.

\subsection{Решаемые задачи}
\begin{enumerate}
  \item \textbf{Поиск.} Сократить время поиска руководителя за счёт фильтрации по темам и ключевым словам вместо ручного перебора списка кафедры.
  \item \textbf{Информированность студента.} Показать статистику отзывов, долю успешных защит и категорию шанса до подачи заявки.
  \item \textbf{Информированность преподавателя.} Дать профиль и оценки студента до личного общения.
  \item \textbf{Коммуникация.} Упростить обмен контактами после заявки.
  \item \textbf{Управление местами.} Учитывать мягкий лимит, избранных и сроки отбора каждого преподавателя.
  \item \textbf{Определённость результата.} Не оставлять студента без руководителя к концу отбора (автоматическое подтверждение по топу).
  \item \textbf{Качество руководства.} Накапливать анонимную оценку руководителей и публичные отзывы с защитой от накруток.
  \item \textbf{Достоверность.} Дать деканату инструмент верификации преподавателей и модерации спорных данных.
\end{enumerate}

\subsection{Ожидаемый результат}
Единая веб-система, где весь цикл «поиск $\rightarrow$ заявка $\rightarrow$ принятие $\rightarrow$ работа $\rightarrow$ отзыв» проходит в одном месте, а решения принимаются на основе данных.

% ================= 3 =================
\section{Требования к системе}

\subsection{Функциональные требования}
{\small
\begin{longtable}{|p{1.4cm}|p{13.4cm}|}
\hline
\textbf{ID} & \textbf{Требование}\\ \hline
\endhead
\multicolumn{2}{|l|}{\textit{Учётные записи и доступ}}\\ \hline
ФТ-1 & Система аутентифицирует пользователей. Целевой способ — вход через ITMO ID (после получения разрешения); в прототипе — локальная аутентификация.\\ \hline
ФТ-2 & Система поддерживает три роли (студент, преподаватель, администратор) с разграничением прав.\\ \hline
ФТ-3 & Преподаватель регистрируется и указывает ФИО, кафедру или направление, контакты (почта, Telegram и др.). Профиль становится виден студентам только после верификации администратором.\\ \hline
ФТ-4 & Система получает данные студента из my.itmo: ФИО, группу, направление, оценки по всем предметам, средний балл, статус и итоговую оценку работы.\\ \hline
\multicolumn{2}{|l|}{\textit{Профиль преподавателя}}\\ \hline
ФТ-5 & Преподаватель создаёт, редактирует и закрывает темы (название, описание, ключевые слова, направление).\\ \hline
ФТ-6 & Преподаватель задаёт лимит мест и может включить запрет приёма сверх лимита. При включённом запрете приём заявок закрывается при достижении лимита.\\ \hline
ФТ-7 & Преподаватель задаёт дату окончания отбора.\\ \hline
\multicolumn{2}{|l|}{\textit{Поиск и профиль}}\\ \hline
ФТ-8 & Студент ищет преподавателей по направлению и ключевым словам (совпадение с темами, в том числе уже сданными) и фильтрует результаты.\\ \hline
ФТ-9 & Профиль преподавателя содержит: темы, агрегированную статистику отзывов в процентах, число завершённых работ, долю успешных защит, текстовые отзывы с именами авторов, число заявок, число избранных, заполненность (принято/лимит), дату окончания отбора.\\ \hline
ФТ-10 & Система показывает студенту категорию шанса (низкий, средний, высокий) для каждого преподавателя.\\ \hline
ФТ-11 & При полной или превышенной заполненности шанс — «низкий», с пояснением, что сверх нормы можно попасть, если понравиться преподавателю (если приём сверх нормы не запрещён).\\ \hline
\multicolumn{2}{|l|}{\textit{Заявки}}\\ \hline
ФТ-12 & Студент подаёт заявку преподавателю, формулируя тему своими словами.\\ \hline
ФТ-13 & Студент имеет не более 5 заявок и упорядочивает их в личный топ. Положение заявки скрыто от преподавателя и не влияет на шанс.\\ \hline
ФТ-14 & Заявка имеет статусы: «На рассмотрении», «Избранное», «Принята», «В работе», «Отклонена», «Отозвана», «Закрыта». Переходы между ними описаны в таблице~\ref{tab:status}.\\ \hline
ФТ-15 & Преподаватель отмечает заявку избранной (и снимает отметку), принимает или отклоняет её. Отклонение требует причины.\\ \hline
ФТ-16 & При просмотре заявки преподаватель видит профиль студента, все его оценки из my.itmo, достижения, средний балл и формулировку темы.\\ \hline
ФТ-17 & Студент подтверждает принятие. При подтверждении все остальные его заявки отзываются.\\ \hline
ФТ-18 & Студент может запросить контакты преподавателя, указанные при регистрации, после подачи заявки.\\ \hline
ФТ-19 & Студент может отозвать свою заявку.\\ \hline
ФТ-20 & Студент может сменить руководителя: система запрашивает причину, закрывает заявку, освобождает место, уведомляет преподавателя и сбрасывает прошлые заявки студента.\\ \hline
\multicolumn{2}{|l|}{\textit{Сроки}}\\ \hline
ФТ-21 & За 24 часа до даты окончания отбора система уведомляет студентов с принятой, но не подтверждённой заявкой.\\ \hline
ФТ-22 & В момент окончания отбора система подтверждает заявку, стоящую в топе студента выше всех остальных, а остальные принятые, но не подтверждённые заявки отзывает.\\ \hline
\multicolumn{2}{|l|}{\textit{Завершение работы и отзывы}}\\ \hline
ФТ-23 & Система получает из my.itmo статус «работа сдана» и итоговую оценку (не сдано — оценка~2) и отражает их в профиле преподавателя.\\ \hline
ФТ-24 & Заявка со сдачей работы закрывается. Студент получает предложение оставить отзыв.\\ \hline
ФТ-25 & Оставить отзыв может только аккаунт со статусом «работа сдана». Оценки по шкалам анонимны (набор шкал уточняется), текстовый комментарий необязателен и публикуется с именем автора.\\ \hline
ФТ-26 & Система пересчитывает агрегированную статистику профиля преподавателя после каждого нового отзыва.\\ \hline
\multicolumn{2}{|l|}{\textit{Администрирование и сервис}}\\ \hline
ФТ-27 & Администратор верифицирует преподавателей, ведёт каталог направлений.\\ \hline
ФТ-28 & Пользователь может подать жалобу на отзыв или спорный статус. Администратор рассматривает её и принимает решение.\\ \hline
ФТ-29 & Система отправляет уведомления об изменении статусов, приближении срока и новых заявках внутри системы (дополнительный канал — почта).\\ \hline
ФТ-30 & Система ведёт журнал значимых событий (смена статусов, решения администратора, ошибки интеграции).\\ \hline
\caption{Функциональные требования}
\end{longtable}}

\subsection{Нефункциональные требования}
{\small
\begin{longtable}{|p{1.4cm}|p{3.2cm}|p{10.2cm}|}
\hline
\textbf{ID} & \textbf{Категория} & \textbf{Требование}\\ \hline
\endhead
НФТ-1 & Производительность & 95\,\% запросов поиска и открытия профиля выполняются не дольше 1\,с при данных порядка 5\,000 студентов и 300 преподавателей (ориентировочные цифры).\\ \hline
НФТ-2 & Нагрузка & Система выдерживает не менее 200 одновременных пользователей, пиковая нагрузка — периоды отбора.\\ \hline
НФТ-3 & Безопасность & Передача по HTTPS, ролевая модель доступа, защита от SQL-инъекций, XSS и CSRF, хранение паролей (в прототипе) только в виде хэшей.\\ \hline
НФТ-4 & Приватность & Оценки студента видны только преподавателям, которым он подал заявку. Автор анонимных оценок не передаётся клиенту ни одним API. Положение заявки в топе недоступно преподавателю на уровне API.\\ \hline
НФТ-5 & Целостность данных & Переходы статусов и подсчёт мест выполняются транзакционно. Обработка сроков идемпотентна, повторный запуск не меняет результат.\\ \hline
НФТ-6 & Надёжность & Недоступность my.itmo не блокирует работу: используются последние загруженные данные, ошибка пишется в журнал.\\ \hline
НФТ-7 & Удобство использования & Интерфейс на русском языке, адаптивная вёрстка. Основные сценарии (поиск, заявка, подтверждение) выполняются не более чем за 3 перехода от главной страницы.\\ \hline
НФТ-8 & Прозрачность & Рядом с категорией шанса отображаются факторы, повлиявшие на неё.\\ \hline
НФТ-9 & Совместимость & Работа в двух последних версиях Chrome, Firefox, Safari и Edge.\\ \hline
НФТ-10 & Сопровождаемость & Слоистая архитектура, версионирование схемы БД миграциями, документация REST API в формате OpenAPI, автоматические тесты бизнес-правил R1–R11.\\ \hline
НФТ-11 & Расширяемость & Интеграция с my.itmo реализуется через интерфейс-адаптер, замена заглушки на реальный клиент не меняет остальной код.\\ \hline
НФТ-12 & Развёртывание & Все этапы работы демонстрируются на сервере \texttt{helios}. Система собирается в одно приложение и запускается без контейнеров и прав суперпользователя.\\ \hline
\caption{Нефункциональные требования}
\end{longtable}}

\subsection{Допущения и ограничения}
\begin{itemize}
  \item Интеграция с my.itmo зависит от разрешения. До его получения данные успеваемости и статус сдачи работы берутся из заглушки.
  \item Набор шкал отзыва и формула расчёта шанса (пороги категорий) уточняются на этапе реализации. Список влияющих факторов зафиксирован в правиле R10.
  \item Автоматического распределения студентов по преподавателям система не выполняет.
\end{itemize}

% ================= 4 =================
\section{Модель прецедентов}

\subsection{Актёры}
\begin{itemize}
  \item \textbf{Студент}, \textbf{Преподаватель}, \textbf{Администратор} — пользователи системы.
  \item \textbf{Планировщик} — системный актёр, запускающий обработку сроков отбора.
  \item \textbf{my.itmo} — внешняя система, поставщик данных об успеваемости и сдаче работ.
\end{itemize}

\subsection{Перечень прецедентов}
\begin{center}
\small
\begin{tabular}{|l|p{8.5cm}|p{3.7cm}|}
\hline
\textbf{ID} & \textbf{Прецедент} & \textbf{Актёр}\\ \hline
UC-1 & Поиск преподавателей & Студент\\ \hline
UC-2 & Просмотр профиля преподавателя и оценка шанса & Студент\\ \hline
UC-3 & Подача заявки и ранжирование топа & Студент\\ \hline
UC-4 & Запрос контактов преподавателя & Студент\\ \hline
UC-5 & Подтверждение выбора преподавателя & Студент\\ \hline
UC-6 & Смена научного руководителя & Студент\\ \hline
UC-7 & Оставление отзыва о преподавателе & Студент\\ \hline
UC-8 & Управление профилем, темами и параметрами отбора & Преподаватель\\ \hline
UC-9 & Рассмотрение заявок & Преподаватель\\ \hline
UC-10 & Обработка окончания срока отбора & Планировщик\\ \hline
UC-11 & Синхронизация данных с my.itmo & my.itmo, Планировщик\\ \hline
UC-12 & Верификация преподавателей & Администратор\\ \hline
UC-13 & Модерация отзывов и жалоб & Администратор\\ \hline
\end{tabular}
\end{center}
Перечень прецедентов предлагается к согласованию с преподавателем.

\subsection{Диаграммы прецедентов}

\begin{figure}[H]
\centering
\resizebox{\textwidth}{!}{%
\begin{tikzpicture}
  \draw[rounded corners=3pt] (3.0,1.2) rectangle (11.0,-10.6);
  \node[font=\small\bfseries] at (7,0.85) {Система поиска научного руководителя};
  \actor{stu}{(0.8,-4.9)}{Студент}
  \node[uc] (u1) at (7,-0.3)  {UC-1. Поиск преподавателей};
  \node[uc] (u2) at (7,-1.85) {UC-2. Просмотр профиля и оценка шанса};
  \node[uc] (u3) at (7,-3.4)  {UC-3. Подача заявки и ранжирование топа};
  \node[uc] (u4) at (7,-4.95) {UC-4. Запрос контактов преподавателя};
  \node[uc] (u5) at (7,-6.5)  {UC-5. Подтверждение выбора};
  \node[uc] (u6) at (7,-8.05) {UC-6. Смена научного руководителя};
  \node[uc] (u7) at (7,-9.6)  {UC-7. Оставление отзыва};
  \foreach \i in {1,...,7} \draw (stu) -- (u\i.west);
\end{tikzpicture}}
\caption{Диаграмма прецедентов. Студент}
\end{figure}

\begin{figure}[H]
\centering
\resizebox{\textwidth}{!}{%
\begin{tikzpicture}
  \draw[rounded corners=3pt] (3.0,1.2) rectangle (11.0,-9.0);
  \node[font=\small\bfseries] at (7,0.85) {Система поиска научного руководителя};
  \actor{tea}{(0.8,-1.0)}{Преподаватель}
  \actor{adm}{(13.0,-4.1)}{Администратор}
  \actor{pla}{(13.0,-6.3)}{Планировщик}
  \actor{itm}{(13.0,-8.3)}{my.itmo}
  \node[uc] (t1) at (7,-0.2)  {UC-8. Управление профилем, темами и параметрами отбора};
  \node[uc] (t2) at (7,-1.75) {UC-9. Рассмотрение заявок};
  \node[uc] (t3) at (7,-3.3)  {UC-12. Верификация преподавателей};
  \node[uc] (t4) at (7,-4.85) {UC-13. Модерация отзывов и жалоб};
  \node[uc] (t5) at (7,-6.4)  {UC-10. Обработка окончания срока отбора};
  \node[uc] (t6) at (7,-7.95) {UC-11. Синхронизация данных с my.itmo};
  \draw (tea) -- (t1.west);
  \draw (tea) -- (t2.west);
  \draw (adm) -- (t3.east);
  \draw (adm) -- (t4.east);
  \draw (pla) -- (t5.east);
  \draw (pla) -- (t6.east);
  \draw (itm) -- (t6.east);
\end{tikzpicture}}
\caption{Диаграмма прецедентов. Преподаватель, администратор и системные актёры}
\end{figure}

\subsection{Модель статусов заявки}
\begin{table}[H]
\centering
\small
\begin{tabular}{|p{3.3cm}|p{2.3cm}|p{3.2cm}|p{5.6cm}|}
\hline
\textbf{Из} & \textbf{В} & \textbf{Инициатор} & \textbf{Условие}\\ \hline
На рассмотрении & Избранное & Преподаватель & Отметка «звёздочкой»; её можно снять (обратный переход)\\ \hline
На рассмотрении, Избранное & Принята & Преподаватель & Отбор открыт; при запрете приёма сверх нормы лимит не исчерпан; у студента нет подтверждённой работы\\ \hline
На рассмотрении, Избранное & Отклонена & Преподаватель & Указана причина\\ \hline
Принята & В работе & Студент или система & Подтверждение студента либо автоподтверждение (UC-10); остальные заявки студента отзываются\\ \hline
На рассмотрении, Избранное, Принята & Отозвана & Студент или система & Отзыв студентом; подтверждение другой заявки; истечение срока без подтверждения (UC-10)\\ \hline
В работе & Закрыта & Студент или система & Смена руководителя (с причиной) либо статус «работа сдана» из my.itmo\\ \hline
\end{tabular}
\caption{Переходы статусов заявки}\label{tab:status}
\end{table}

\subsection{Описание прецедентов}

\begin{ucd}{UC-1}{Поиск преподавателей}{Студент}
Предусловия & Студент аутентифицирован.\\ \hline
Основной сценарий & \steps{
 \item Студент указывает направление и ключевые слова или темы интереса.
 \item Система сопоставляет запрос с темами преподавателей (текущими и сданными) и ключевыми словами. В выдачу попадают только верифицированные преподаватели.
 \item Система выводит список: ФИО, кафедра, совпавшие темы, категория шанса, признак «отбор открыт» и дата его окончания.
 \item Студент фильтрует и сортирует список и открывает нужный профиль (UC-2).}\\ \hline
Альтернативы & 2а. Совпадений нет: система предлагает ослабить условия и показывает преподавателей выбранного направления.\\ \hline
Постусловия & Студенту показан список подходящих преподавателей.\\ \hline
\end{ucd}

\begin{ucd}{UC-2}{Просмотр профиля преподавателя и оценка шанса}{Студент}
Предусловия & Студент выбрал преподавателя из списка.\\ \hline
Основной сценарий & \steps{
 \item Система показывает темы (текущие и сданные), статистику отзывов в процентах, число завершённых работ, долю успешных защит и текстовые отзывы с именами авторов.
 \item Система показывает число заявок, число избранных («почти принятых»), заполненность (принято/лимит) и дату окончания отбора.
 \item Система рассчитывает и показывает категорию шанса для этого студента с перечнем факторов (правило R10).}\\ \hline
Альтернативы & 3а. Лимит исчерпан или превышен: шанс «низкий» с пояснением о возможности попасть сверх нормы.\newline
3б. Включён запрет приёма сверх нормы и лимит достигнут: приём заявок закрыт, подача недоступна.\\ \hline
Постусловия & Студент владеет данными для решения о подаче заявки.\\ \hline
\end{ucd}

\begin{ucd}{UC-3}{Подача заявки и ранжирование топа}{Студент}
Предусловия & Нет подтверждённой работы; активных заявок меньше пяти; отбор у преподавателя открыт.\\ \hline
Основной сценарий & \steps{
 \item Студент нажимает «Подать заявку» в профиле преподавателя.
 \item Студент формулирует тему своими словами.
 \item Система создаёт заявку со статусом «На рассмотрении» и ставит её в конец топа студента.
 \item Система уведомляет преподавателя.
 \item Студент изменяет порядок заявок в топе. Рядом с каждой показана категория шанса. Топ видит только студент.}\\ \hline
Альтернативы & 1а. Уже пять активных заявок: подача отклоняется, предлагается отозвать одну из них.\newline
5а. Студент отзывает заявку: статус «Отозвана», место в топе освобождается.\\ \hline
Постусловия & Заявка зарегистрирована. Положение в топе скрыто от преподавателя.\\ \hline
\end{ucd}

\begin{ucd}{UC-4}{Запрос контактов преподавателя}{Студент}
Предусловия & У студента есть заявка этому преподавателю, не отклонённая и не отозванная.\\ \hline
Основной сценарий & \steps{
 \item Студент выбирает «Показать контакты» в карточке заявки.
 \item Система показывает контакты, указанные преподавателем при регистрации (почта, Telegram и др.).
 \item Дальнейшее общение проходит вне системы.}\\ \hline
Альтернативы & 2а. Контакты не заполнены: система сообщает об этом и предлагает дождаться ответа преподавателя.\\ \hline
Постусловия & Студент получил способ связаться с преподавателем.\\ \hline
\end{ucd}

\begin{ucd}{UC-5}{Подтверждение выбора преподавателя}{Студент}
Предусловия & Преподаватель перевёл заявку в статус «Принята».\\ \hline
Основной сценарий & \steps{
 \item Система уведомляет студента о принятии.
 \item Студент открывает заявку и подтверждает, что идёт к этому преподавателю.
 \item Система переводит заявку в «В работе» и создаёт запись о работе (студент, преподаватель, тема).
 \item Система отзывает все остальные заявки студента и очищает его топ.
 \item Система уведомляет преподавателей, чьи заявки отозваны.}\\ \hline
Альтернативы & 2а. Студент отказывается от принятой заявки: статус «Отозвана», остальные заявки сохраняются.\newline
2б. Студент не реагирует до даты окончания отбора: выполняется UC-10.\\ \hline
Постусловия & Студент закреплён за преподавателем, место в лимите занято.\\ \hline
\end{ucd}

\begin{ucd}{UC-6}{Смена научного руководителя}{Студент}
Предусловия & Заявка студента в статусе «В работе».\\ \hline
Основной сценарий & \steps{
 \item Студент выбирает «Сменить руководителя».
 \item Система запрашивает причину.
 \item Студент указывает причину.
 \item Система закрывает заявку, сохраняет причину, освобождает место в лимите и уведомляет преподавателя.
 \item Система сбрасывает все прошлые заявки студента.
 \item Студент возвращается к поиску (UC-1) и подаёт новые заявки (UC-3).}\\ \hline
Постусловия & Заявка закрыта, место освобождено, у студента нет активных заявок.\\ \hline
\end{ucd}

\begin{ucd}{UC-7}{Оставление отзыва о преподавателе}{Студент}
Предусловия & Заявка получила статус «работа сдана» (UC-11); отзыв по ней ещё не оставлен.\\ \hline
Основной сценарий & \steps{
 \item Система предлагает студенту оставить отзыв.
 \item Студент выставляет оценки по шкалам (строгость, гибкость сроков, вовлечённость и др.). Оценки анонимны.
 \item Студент по желанию вводит текстовый комментарий, который публикуется с его именем.
 \item Система сохраняет отзыв так, чтобы автор оценок не передавался клиентам.
 \item Система пересчитывает статистику профиля преподавателя. Преподаватель видит обновлённую статистику и комментарий.}\\ \hline
Альтернативы & 2а. Студент не оставляет отзыв: ничего не происходит, система может напомнить позже.\\ \hline
Постусловия & Статистика преподавателя обновлена.\\ \hline
\end{ucd}

\begin{ucd}{UC-8}{Управление профилем, темами и параметрами отбора}{Преподаватель}
Предусловия & Преподаватель зарегистрирован. Для публикации профиля требуется верификация (UC-12).\\ \hline
Основной сценарий & \steps{
 \item Преподаватель заполняет профиль: кафедра, контакты.
 \item Преподаватель добавляет и редактирует темы (название, описание, ключевые слова, направление) и закрывает неактуальные.
 \item Преподаватель задаёт лимит мест и при необходимости включает запрет приёма сверх нормы.
 \item Преподаватель задаёт дату окончания отбора.
 \item Система сохраняет изменения и пересчитывает категории шансов.}\\ \hline
Альтернативы & 0а. Аккаунт не верифицирован: профиль не виден студентам.\newline
3а. Лимит уменьшен ниже числа принятых: действующие студенты сохраняются, приём новых закрывается.\\ \hline
Постусловия & Профиль актуален и виден студентам.\\ \hline
\end{ucd}

\begin{ucd}{UC-9}{Рассмотрение заявок}{Преподаватель}
Предусловия & У преподавателя есть входящие заявки.\\ \hline
Основной сценарий & \steps{
 \item Преподаватель открывает список входящих заявок.
 \item Преподаватель открывает заявку: профиль студента, все оценки из my.itmo, средний балл, достижения, формулировка темы.
 \item Преподаватель выбирает решение: отметить избранным, принять или отклонить.
 \item Система меняет статус и уведомляет студента.}\\ \hline
Альтернативы & 3а. Отклонение: причина обязательна.\newline
3б. Принятие при включённом запрете приёма сверх нормы и исчерпанном лимите блокируется. Без запрета принятие разрешено, заполненность может превысить 100\,\%.\newline
3в. Студент уже в статусе «В работе» у другого преподавателя: принятие недоступно.\\ \hline
Постусловия & Заявка в новом статусе, студент уведомлён.\\ \hline
\end{ucd}

\begin{ucd}{UC-10}{Обработка окончания срока отбора}{Планировщик}
Предусловия & У преподавателя задана дата окончания отбора.\\ \hline
Основной сценарий & \steps{
 \item За 24 часа до даты система отправляет напоминание каждому студенту с принятой, но не подтверждённой заявкой.
 \item В момент окончания для каждой такой заявки система проверяет её положение в топе студента.
 \item Если заявка выше всех остальных заявок студента, она подтверждается автоматически (как в UC-5, шаги 3--5).
 \item Если выше стоит другая заявка, текущая отзывается.}\\ \hline
Альтернативы & 2а. Студент подтвердил или отозвал заявку до срока: заявка не обрабатывается.\\ \hline
Постусловия & Принятые заявки подтверждены или отозваны, повторный запуск ничего не меняет.\\ \hline
\end{ucd}

\begin{ucd}{UC-11}{Синхронизация данных с my.itmo}{my.itmo, Планировщик}
Предусловия & Настроен доступ к my.itmo (до получения разрешения — заглушка).\\ \hline
Основной сценарий & \steps{
 \item Система запрашивает оценки и средний балл студентов с активными заявками.
 \item Система запрашивает статус и итоговую оценку работ в статусе «В работе».
 \item Для сданных работ система фиксирует «работа сдана» и оценку (не сдано — 2), закрывает заявку и открывает студенту возможность оставить отзыв (UC-7).
 \item Система обновляет статистику в профиле преподавателя.}\\ \hline
Альтернативы & 1а. my.itmo недоступна: система использует последние данные и записывает ошибку в журнал.\\ \hline
Постусловия & Данные об успеваемости и сдаче актуальны.\\ \hline
\end{ucd}

\begin{ucd}{UC-12}{Верификация преподавателей}{Администратор}
Предусловия & Есть неподтверждённые аккаунты преподавателей.\\ \hline
Основной сценарий & \steps{
 \item Администратор открывает очередь на верификацию.
 \item Администратор проверяет данные преподавателя.
 \item Администратор подтверждает или отклоняет аккаунт.
 \item Система публикует профиль или оставляет его скрытым и уведомляет преподавателя.}\\ \hline
Постусловия & Каталог содержит только подтверждённых преподавателей.\\ \hline
\end{ucd}

\begin{ucd}{UC-13}{Модерация отзывов и жалоб}{Администратор}
Предусловия & Поступила жалоба на отзыв или спорный статус заявки (ФТ-28).\\ \hline
Основной сценарий & \steps{
 \item Администратор открывает жалобу и связанные данные.
 \item Администратор принимает решение: оставить отзыв, скрыть его или исправить спорный статус.
 \item Система применяет решение, пересчитывает статистику при необходимости и уведомляет заявителя.}\\ \hline
Постусловия & Спорный вопрос закрыт, решение записано в журнал.\\ \hline
\end{ucd}

% ================= 5 =================
\section{Архитектура системы}

\subsection{Общее описание}
Система реализуется как клиент-серверное веб-приложение. Одностраничный клиент (SPA) взаимодействует с серверным приложением по REST/JSON поверх HTTPS. Сервер состоит из слоёв: веб-слой, сервисный слой, слой доступа к данным. Рядом с сервисным слоем расположены планировщик (сроки отбора, синхронизация) и адаптер внешней системы my.itmo. Данные хранятся в PostgreSQL. Серверное приложение и база данных размещаются на сервере \texttt{helios}.

\begin{figure}[H]
\centering
\resizebox{\textwidth}{!}{%
\begin{tikzpicture}
  \draw[dashed,rounded corners=3pt] (3.2,0.4) rectangle (13.0,-8.9);
  \node[anchor=north west,font=\small\bfseries] at (3.3,0.35) {Сервер helios};
  \draw[rounded corners=3pt] (3.6,-0.5) rectangle (12.4,-6.6);
  \node[anchor=north west,font=\small] at (3.7,-0.55) {Серверное приложение (Java, Spring Boot)};

  \node[box,minimum width=2.9cm,minimum height=1.4cm] (br) at (-1.4,-1.6) {Браузер\\Vue SPA};
  \node[box,minimum width=7.6cm] (web) at (8.0,-1.6) {Веб-слой: REST-контроллеры, Spring Security};
  \node[box,minimum width=7.6cm,minimum height=1.2cm] (svc) at (8.0,-3.1) {Сервисный слой\\поиск · заявки · расчёт шанса · отзывы · уведомления};
  \node[box,minimum width=3.5cm] (sch) at (5.9,-4.6) {Планировщик\\(сроки, синхронизация)};
  \node[box,minimum width=3.5cm] (ada) at (10.3,-4.6) {Адаптер my.itmo};
  \node[box,minimum width=7.6cm] (dal) at (8.0,-5.9) {Доступ к данным: Spring Data JPA, миграции};
  \node[box,minimum width=3.5cm,minimum height=1.0cm] (db) at (8.0,-7.9) {PostgreSQL};
  \node[box,minimum width=3.0cm,minimum height=1.2cm] (ext) at (15.6,-4.6) {my.itmo\\(внешняя система)};

  \draw[<->] (br.east) -- node[above,font=\scriptsize,align=center,fill=white,inner sep=1pt] {HTTPS\\REST/JSON} (web.west);
  \draw[<->] (web) -- (svc);
  \draw[->] (sch.north) -- (sch.north |- svc.south);
  \draw[->] (ada.north) -- (ada.north |- svc.south);
  \draw[<->] (svc.south) -- (svc.south |- dal.north);
  \draw[<->] (dal) -- (db);
  \draw[<->] (ada.east) -- (ext.west);
\end{tikzpicture}}
\caption{Архитектура и размещение компонентов}
\end{figure}

\subsection{Компоненты серверного приложения}
\begin{center}
\small
\begin{tabular}{|p{3.7cm}|p{11.1cm}|}
\hline
\textbf{Компонент} & \textbf{Ответственность}\\ \hline
Аутентификация и роли & Вход пользователей, роли студент/преподаватель/администратор, контроль доступа к методам API\\ \hline
Каталог и поиск & Профили преподавателей, темы, поиск по ключевым словам и направлению, верификация\\ \hline
Заявки & Создание, топ студента, статусы и переходы из таблицы~\ref{tab:status}, проверка правил R1–R8\\ \hline
Расчёт шанса & Вычисление категории по правилу R10 и перечня факторов для пояснения\\ \hline
Отзывы и работы & Фиксация сдачи, анонимные оценки, публичные комментарии, агрегаты, жалобы\\ \hline
Уведомления & Внутренние уведомления, опциональная отправка по почте\\ \hline
Планировщик & Напоминание за сутки, обработка окончания отбора (UC-10), периодическая синхронизация (UC-11)\\ \hline
Адаптер my.itmo & Интерфейс получения оценок и статуса работы. Две реализации: заглушка с тестовыми данными и HTTP-клиент\\ \hline
\end{tabular}
\end{center}

\subsection{Модель данных}
Основные сущности (уточнённая версия модели из предыдущего задания):
\begin{center}
\small
\begin{tabular}{|p{3cm}|p{11.8cm}|}
\hline
\textbf{Сущность} & \textbf{Основные поля}\\ \hline
Пользователь & id, роль, ФИО, логин, почта\\ \hline
Студент & id, группа, курс, направление, уровень (бакалавр/магистр), достижения; оценки и средний балл подтягиваются из my.itmo\\ \hline
Преподаватель & id, кафедра, контакты, лимит мест, запрет приёма сверх нормы, дата окончания отбора, статус верификации\\ \hline
Тема & id, преподаватель, название, описание, ключевые слова, направление, статус (открыта/закрыта), признак «сдана»\\ \hline
Заявка & id, студент, преподаватель, формулировка темы, статус, позиция в топе студента, избранное, причина отклонения или закрытия, даты\\ \hline
Работа & id, заявка, статус (в работе/сдана), итоговая оценка\\ \hline
Отзыв & id, работа, оценки по шкалам, текст (необязательно), дата. Автор хранится отдельно и не отдаётся через API\\ \hline
Жалоба & id, автор, объект (отзыв или заявка), причина, решение\\ \hline
Уведомление & id, получатель, тип, текст, прочитано\\ \hline
\end{tabular}
\end{center}
Связи: преподаватель владеет темами (1:N); студент и преподаватель связаны через заявку (M:N); работа порождается подтверждённой заявкой (1:1); отзыв принадлежит работе (1:1).

\begin{figure}[H]
\centering
\resizebox{\textwidth}{!}{%
\begin{tikzpicture}[
  ent/.style={draw,rectangle,align=left,font=\footnotesize,inner sep=3pt,fill=white},
  card/.style={fill=white,inner sep=1pt,font=\scriptsize},
  rel/.style={fill=white,inner sep=1pt,font=\scriptsize\itshape}
]
% --- сущности ---
\node[ent] (user) at (7,7.5) {\ent{ПОЛЬЗОВАТЕЛЬ}{\underline{id}\\ role\\ full\_name\\ login\\ email}};
\node[ent] (notif) at (0,7.5) {\ent{УВЕДОМЛЕНИЕ}{\underline{id}\\ user\_id\ (FK)\\ type\\ text\\ is\_read\\ created\_at}};
\node[ent] (stud) at (0,0) {\ent{СТУДЕНТ}{\underline{user\_id}\ (FK)\\ study\_group\\ course\\ direction\\ level\\ achievements}};
\node[ent] (appl) at (7,0) {\ent{ЗАЯВКА}{\underline{id}\\ student\_id\ (FK)\\ teacher\_id\ (FK)\\ topic\_text\\ status\\ rank\\ is\_favorite\\ reason\\ created\_at}};
\node[ent] (teach) at (14,0) {\ent{ПРЕПОДАВАТЕЛЬ}{\underline{user\_id}\ (FK)\\ department\\ contacts\\ quota\\ no\_over\_quota\\ selection\_end\\ verified}};
\node[ent] (topic) at (14,-6) {\ent{ТЕМА}{\underline{id}\\ teacher\_id\ (FK)\\ title\\ description\\ keywords\\ direction\\ is\_open\\ is\_completed}};
\node[ent] (work) at (7,-5.5) {\ent{РАБОТА}{\underline{id}\\ application\_id\ (FK)\\ status\\ grade}};
\node[ent] (rev) at (7,-12) {\ent{ОТЗЫВ}{\underline{id}\\ work\_id\ (FK)\\ strictness\\ flexibility\\ engagement\\ comment\\ created\_at}};
\node[ent] (rauth) at (0,-9.8) {\ent{АВТОР ОТЗЫВА}{\underline{review\_id}\ (FK)\\ student\_id\ (FK)}};

% --- связи ---
\draw (user) -- (stud)
  node[card,pos=0.1]{1} node[card,pos=0.9]{0..1}
  node[rel,pos=0.5]{роль студента};
\draw (user) -- (teach)
  node[card,pos=0.1]{1} node[card,pos=0.9]{0..1}
  node[rel,pos=0.5]{роль преподавателя};
\draw (user) -- (notif)
  node[card,pos=0.12]{1} node[card,pos=0.88]{N}
  node[rel,pos=0.5]{получает};
\draw (stud) -- (appl)
  node[card,pos=0.12]{1} node[card,pos=0.88]{N}
  node[rel,pos=0.5]{подаёт};
\draw (teach) -- (appl)
  node[card,pos=0.12]{1} node[card,pos=0.88]{N}
  node[rel,pos=0.5]{получает};
\draw (teach) -- (topic)
  node[card,pos=0.1]{1} node[card,pos=0.9]{N}
  node[rel,pos=0.5]{ведёт};
\draw (appl) -- (work)
  node[card,pos=0.1]{1} node[card,pos=0.9]{0..1}
  node[rel,pos=0.5]{порождает};
\draw (work) -- (rev)
  node[card,pos=0.1]{1} node[card,pos=0.9]{0..1}
  node[rel,pos=0.5]{получает};
\draw (rev) -- (rauth)
  node[card,pos=0.1]{1} node[card,pos=0.9]{1}
  node[rel,pos=0.5]{автор};
\end{tikzpicture}}
\caption{ER-модель системы}
\end{figure}

\subsection{Технологии}
Выбор соответствует разрешённым для курсовой технологиям и подлежит согласованию с преподавателем.
\begin{center}
\small
\begin{tabular}{|p{2.5cm}|p{6.2cm}|p{5.9cm}|}
\hline
\textbf{Слой} & \textbf{Технология} & \textbf{Обоснование}\\ \hline
Frontend & Vue 3, TypeScript, Vite, Vue Router, Pinia & Интерактивные фильтры, перетаскивание заявок в топе, низкий порог входа\\ \hline
Backend & Java, Spring Boot (Spring MVC), Spring Security, Spring Data JPA, Spring Scheduling, Bean Validation & Готовая поддержка REST, ролей и планировщика\\ \hline
БД & PostgreSQL, миграции Flyway & Требуется курсом, транзакционность нужна для счёта мест и переходов статусов\\ \hline
API & REST/JSON, OpenAPI & Чёткий контракт между клиентом и сервером\\ \hline
Сборка & Maven (или Gradle), npm & Фронтенд собирается в статические файлы, которые раздаёт приложение\\ \hline
Тесты & JUnit 5, Mockito & Автоматическая проверка правил R1–R11\\ \hline
\end{tabular}
\end{center}
Альтернативой клиентской части может быть Spring MVC с Thymeleaf, что упростит развёртывание, но ограничит интерактивность.

\subsection{Размещение на сервере \texttt{helios}}
\begin{itemize}
  \item Клиентская часть собирается в статические файлы и входит в состав исполняемого JAR серверного приложения, поэтому на \texttt{helios} запускается одно приложение.
  \item PostgreSQL используется как СУБД, развёрнутая на том же сервере. Параметры подключения вынесены в конфигурацию.
  \item Каждый этап курсовой собирается и запускается на \texttt{helios} без контейнеров и прав суперпользователя.
\end{itemize}


\end{document}